% Same-visit check on the whole sample (input from cr_pair_results.tex)
\subsection{The whole sample}
\label{sec:cr_pair_sample}
We reduced all 106 exposures of the sample again with the check (60 of NGC\,4151 at OGG, 46 of Fairall\,9 at COJ).
Every one found its partner (\texttt{L1CRMETH} = \texttt{pair}). These products were made with the version of the
check \textit{before} the rule for pixels flagged in both exposures on the object (Section~\ref{sec:cr_pair}, step 8).
That rule was added after this run and measured on the seven frames with the most flags left on the lines; the sample
was not reduced again. Below, we call the version run on the sample the \textit{sample version}, and each number says
which version it comes from.

\paragraph{Flags on the narrow lines.} Table~\ref{tab:cr_pair_lines} counts the flagged pixels within 12\,\AA\ of
[O\,\textsc{iii}]\,$\lambda\lambda$4959,\,5007, H$\beta$ and H$\alpha$+[N\,\textsc{ii}] and within 8\,pixels of the
trace. On the NGC\,4151 lines, the first version of the check doubled the flags of the object model. The sample
version brings them to 215, fewer than with the object model alone (398) and with 1.1.0 (236), and it clears most of
the H$\alpha$ flags of 1.1.0. Fairall\,9 has few flags on its lines with any detector.

Most of these 215 pixels are still not cosmic rays. They are 0.32 per cent of the pixels in the line windows, against
0.003 per cent in five continuum windows on the same traces. Often both exposures of a visit are flagged at the same
line, which a cosmic ray cannot do. We found two causes in the worst frames. First, in good-seeing visits LA Cosmic
flags the same line and trace peaks in both exposures, so neither can check the other (item (v) of
Section~\ref{sec:cr_pair}). Second, on tilted lines a residual offset between the exposures leaves a
positive--negative pattern, whose positive half passes as a confirmed cosmic ray. The rule for pixels flagged in both
exposures removes the first cause (Table~\ref{tab:cr_pair_worst}). We left the second, because a rule for it would be
one more cut tuned on the frames it is meant to fix.

\begin{table}[!htbp]
\centering\small
\caption{Pixels flagged as cosmic rays on the narrow lines of the sample (within 12\,\AA\ of the line and 8\,pixels of
the trace), by order (1 red, 2 blue), for the 1.1.0 detector (improved v3 products), LA Cosmic with the object model, the
first version of the same-visit check, and the version run on the sample (before the rule for pixels flagged in both
exposures). Last rows: total, and frames with any flagged pixel.}
\label{tab:cr_pair_lines}
\begin{tabular}{llcccc}
\toprule
Target & Line, order & 1.1.0 & Object model & Check, first & Check, sample \\
 & & \multicolumn{4}{c}{(pixels)} \\
\midrule
NGC\,4151 & H$\alpha$+[N\,\textsc{ii}], 1 & 211 & 70 & 55 & 20 \\
(60 frames) & H$\beta$, 1 & 1 & 56 & 117 & 28 \\
 & H$\beta$, 2 & 0 & 6 & 0 & 0 \\
 & [O\,\textsc{iii}]\,4959, 1 & 0 & 42 & 71 & 31 \\
 & [O\,\textsc{iii}]\,4959, 2 & 24 & 74 & 139 & 45 \\
 & [O\,\textsc{iii}]\,5007, 1 & 0 & 89 & 170 & 77 \\
 & [O\,\textsc{iii}]\,5007, 2 & 0 & 61 & 288 & 14 \\
 & total & 236 & 398 & 840 & 215 \\
 & frames with any & 18 & 37 & 41 & 28 \\
\midrule
Fairall\,9 & total & 12 & 29 & 21 & 15 \\
(46 frames) & frames with any & 2 & 6 & 4 & 4 \\
\bottomrule
\end{tabular}
\end{table}

\begin{table}[!htbp]
\centering\small
\caption{The rule for pixels flagged in both exposures on the seven NGC\,4151 frames with the most flags on the
lines: pixels on the lines (as in Table~\ref{tab:cr_pair_lines}) flagged by LA Cosmic with the object model, by the
version run on the sample, and with the rule; and the LA Cosmic groups kept because the partner cannot check them
(mostly warm pixels, which the rule must not touch).}
\label{tab:cr_pair_worst}
\begin{tabular}{lccccc}
\toprule
Frame & \multicolumn{3}{c}{Pixels on the lines} & \multicolumn{2}{c}{Groups the partner cannot check} \\
 & LA + object model & sample version & with the rule & sample version & with the rule \\
\midrule
2025-03-27 0021 & 32 & 33 & 20 & 56 & 47 \\
2025-03-27 0020 & 24 & 19 & 13 & 56 & 48 \\
2025-05-31 0020 & 18 & 15 & 12 & 35 & 32 \\
2025-06-10 0006 & 21 & 26 & 22 & 6 & 5 \\
2025-06-10 0007 & 18 & 16 & 16 & 5 & 5 \\
2025-07-05 0004 & 17 & 14 & 14 & 23 & 23 \\
2025-07-14 0002 & 23 & 0 & 0 & 34 & 30 \\
\bottomrule
\end{tabular}
\end{table}

The rule is restricted to the object for a reason. When we applied it everywhere, it took the groups the partner
cannot check from 56 to 0 in the first frame of Table~\ref{tab:cr_pair_worst}. Cutouts showed that the extra groups it
removed off the trace were single warm pixels present in both exposures. On the object, the rule removes only line and
trace peaks, but it will also unflag a warm pixel that happens to lie on the trace. From the density of warm pixels
flagged in both exposures on the sky, we expect 4--7 such groups per frame. Each adds 50--90\,e$^-$ (at most
340\,e$^-$) to a column of the trace, which holds 30\,000--70\,000\,e$^-$ (6000\,e$^-$ in its faintest tenth): typically 0.1--0.3 per
cent of one of the 1700 columns, and at most $\approx6$ per cent (340\,e$^-$ in a column of 6000\,e$^-$). A map of the warm pixels, made from the flags that recur at the same place on the
detector in many frames, would remove them properly; we leave it for later.

\paragraph{The header keywords.} The keywords of the check give, per frame (median, with the range in brackets):
\begin{itemize}
\item \textit{OGG}: LA Cosmic flags 242 (115--715) pixels, of which 188 (96--681) remain after the check (80 per cent,
summed over all frames). 9 groups are confirmed, 12 removed, 34 kept because the partner cannot check them, and 0 (at
most 6) added. The partner cannot judge 0.1 per cent of the orders (at most 6 per cent). The offset along the
dispersion is at most 0.6\,pixel.
\item \textit{COJ}: 321 (175--1755) pixels flagged by LA Cosmic become 278 (129--1519), 90 per cent. 14 groups are
confirmed, 11 removed, 47 kept and 1 (at most 3) added. The partner cannot judge 0.2 per cent of the orders (at most
2.5 per cent).
\end{itemize}

\paragraph{Two faults found on the sample, and fixed.} The sample numbers in this section were measured before these
fixes.
\begin{enumerate}
\item[(i)] \textit{Spurious offsets along the dispersion.} In the blue order of six Fairall\,9 frames, the offset of
the partner along the dispersion came out at 1.4--3.4\,pixels, with the same sign for both exposures of a visit (2026
October 26: $-2.05$ and $-2.05$\,pixels; September 30: $-3.37$ and $-1.52$). A real offset cannot do this: if $b$ is
shifted with respect to $a$, then $a$ is shifted by the opposite amount with respect to $b$. The faint blue spectrum of
Fairall\,9 has no features above the noise. There the normalised cross-correlation peaks at only 0.09--0.26, against
0.92--1.00 in its red order and in both orders of NGC\,4151, and its peak lands wherever the noise puts it: leaving the
flagged cosmic-ray pixels out of the correlation moved one offset from $-2.05$ to $+2.22$\,pixels. An offset is now
used only where the peak exceeds 0.5, which separates the two groups cleanly.
\item[(ii)] \textit{New groups at the end of an order.} In the same frames the check added groups of 1--9 pixels right
at the end of the blue order (1\,pixel from it). There the signal stops sharply, and the matched partner misses the
step by a pixel. The 6-pixel margin in which no new group is added covered the edges of the orders along the slit, but
not their ends along the dispersion. It now covers the ends too.
\end{enumerate}
Measured again on the affected frames, the spurious offsets are gone (the red order keeps its offsets of
0.05--0.24\,pixel), and no group is added at an order end. The flags on the lines of the NGC\,4151 frames of
Table~\ref{tab:cr_pair_worst} are unchanged. Both faults have their own unit test (16 in all).

\paragraph{Fluxes.} What matters in the end is the flux. Table~\ref{tab:cr_pair_flux} gives the robust scatter of the
flux ratio of the two exposures of each visit, for the narrow lines and two continuum windows. This scatter is set by
the slit losses and the seeing, and the choice of detector changes it by 0.5 per cent or less. The one larger change,
for the red-order [O\,\textsc{iii}] of NGC\,4151 with the object model alone (2.5 against 1.7 per cent), is mostly
undone by the check (2.0 per cent). The optimal extraction rebuilds a masked pixel from the profile, so a false flag on
a line peak costs little flux, while a missed cosmic ray costs more. The gain that matters for the light curves is
therefore that of the object model: of the hits above 10$\sigma$ on the OGG line peaks, it finds 89 per cent, where
1.1.0 found 27 per cent (Table~\ref{tab:cr_pair_inject}). The same-visit check is a clean-up on top: 10--22 per cent
fewer flagged pixels (22 at OGG, 10 at COJ), a clean H$\alpha$ peak, and a mask that can be trusted on the lines.

\begin{table}[!htbp]
\centering\small
\caption{Robust scatter of the flux ratio of the two exposures of each visit (30 visits of NGC\,4151 at OGG, 23 of
Fairall\,9 at COJ), for the 1.1.0 detector (improved v3 products), LA Cosmic with the object model, and the same-visit check
(version run on the sample). [O\,\textsc{iii}]\,5007 from a doublet fit in each order; H$\alpha$+[N\,\textsc{ii}] for
NGC\,4151 only.}
\label{tab:cr_pair_flux}
\begin{tabular}{llccc}
\toprule
Target & Measurement & 1.1.0 & Object model & Same-visit check \\
 & & \multicolumn{3}{c}{(scatter of the ratio, per cent)} \\
\midrule
NGC\,4151 & [O\,\textsc{iii}]\,5007, blue order & 2.0 & 1.9 & 2.0 \\
 & [O\,\textsc{iii}]\,5007, red order & 1.7 & 2.5 & 2.0 \\
 & H$\alpha$+[N\,\textsc{ii}] & 2.5 & 2.3 & 2.5 \\
 & continuum, blue & 1.5 & 1.6 & 1.5 \\
 & continuum, red & 1.5 & 1.6 & 1.5 \\
\midrule
Fairall\,9 & [O\,\textsc{iii}]\,5007, blue order & 6.7 & 6.8 & 6.4 \\
 & [O\,\textsc{iii}]\,5007, red order & 6.7 & 6.6 & 6.8 \\
 & continuum, blue & 3.3 & 3.6 & 3.6 \\
 & continuum, red & 3.0 & 2.8 & 2.8 \\
\bottomrule
\end{tabular}
\end{table}
